\section{Copy-Block Harmful Excess And Residue Energy}
\label{sec:copy-block-harmful-excess}

The harmful excess of a copy block is not an arbitrary scalar. It is exactly
the sum of two centered entries of the old start histogram modulo $r$, for
every old period $M\ge1$, every incoming prime $r\ge5$ coprime to $M$, every
run of complete copy blocks, and every finite integer interval in the copied
stream, over lifted copies of one complete old-period 2-gap-start set,
grouped into old-period copy blocks. This turns residue-collision energy into
a quantitative bound for the complete block portion of a local interval.

Let $S\subset[0,M)$ be the old 2-gap starts and put $N=|S|$. For
$t\bmod r$, define
\begin{equation*}
c_t=\#\{a\in S:a\equiv t\pmod r\},
\qquad
d_t=c_t-\frac Nr,
\qquad
V_r=\sum_{t\bmod r}d_t^2.
\end{equation*}

Copy block $j$ is $[jM,(j+1)M)$. Let $K_j$ count starts $a+jM$ destroyed by
filter $r$ and define its centered harmful excess
\begin{equation*}
B_j=K_j-\frac{2N}{r}.
\end{equation*}

Set $t_j\equiv-jM\pmod r$. The two endpoint strikes are disjoint because
$r>2$, and they are equivalent to
\begin{equation*}
a\equiv t_j\pmod r,
\qquad
a\equiv t_j-2\pmod r.
\end{equation*}

Consequently,
\begin{equation*}
\begin{aligned}
K_j
&=c_{t_j}+c_{t_j-2}
&&[\text{Two Endpoint Classes}]\\
B_j
&=d_{t_j}+d_{t_j-2}.
&&\blacksquare\ \text{[Centering; Q.E.D.]}
\end{aligned}
\end{equation*}

Because $\gcd(M,r)=1$, the map $j\mapsto-jM\pmod r$ permutes all residues.
Also $\sum_td_t=\sum_tc_t-N=0$. Hence a complete run of $r$ blocks has
\begin{equation*}
\begin{aligned}
\sum_{j=0}^{r-1}B_j
&=\sum_{t\bmod r}(d_t+d_{t-2})
&&[\text{Residue Permutation}]\\
&=2\sum_{t\bmod r}d_t
&&[\text{Cyclic Reindexing}]\\
&=0.
&&[\text{Centered Histogram}]
\end{aligned}
\end{equation*}

The same permutation gives the exact energy identity
\begin{equation*}
\begin{aligned}
\sum_{j=0}^{r-1}B_j^2
&=\sum_{t\bmod r}(d_t+d_{t-2})^2
&&[\text{Residue Permutation}]\\
&=2V_r+2\sum_{t\bmod r}d_td_{t-2}.
&&[\text{Expansion}]
\end{aligned}
\end{equation*}

The autocorrelation may be negative. Discarding its sign with
$2xy\le x^2+y^2$ gives
\begin{equation*}
\begin{aligned}
2\sum_td_td_{t-2}
&\le\sum_td_t^2+\sum_td_{t-2}^2
&&[\text{Termwise Quadratic Bound}]\\
&=2V_r,
&&[\text{Cyclic Reindexing}]\\
\sum_{j=0}^{r-1}B_j^2
&\le4V_r.
&&\blacksquare\ \text{[Substitution; Q.E.D.]}
\end{aligned}
\end{equation*}

For any $k$ consecutive blocks with $0\le k\lt r$, Cauchy--Schwarz yields
\begin{equation*}
\begin{aligned}
\left|\sum_{j\in J}B_j\right|^2
&\le k\sum_{j\in J}B_j^2
&&[\text{Cauchy--Schwarz}]\\
&\le k\sum_{j=0}^{r-1}B_j^2
&&[\text{Nonnegative Terms}]\\
&\le4kV_r.
&&[\text{Complete-Block Energy}]
\end{aligned}
\end{equation*}

Thus $|\sum_{j\in J}B_j|\le2\sqrt{kV_r}$. For a longer run, remove complete
groups of $r$ blocks by the zero-sum identity and take $k$ to be the number of
remaining blocks.

An arbitrary integer interval $I$ has a left partial old-period block, a run
of complete blocks, and a right partial block. Each partial block contains at
most one copy of every start in $S$. Since $r\ge5$ and the strike indicator is
either $0$ or $1$,
\begin{equation*}
\left|
\mathbf1_{r\mid x(x+2)}-\frac2r
\right|
\le1-\frac2r.
\end{equation*}

Each partial contribution is therefore at most $N(1-2/r)$ in absolute value.
Combining both fragments with the complete-block bound gives
\begin{equation*}
|b_r(I)|
\le
2N\left(1-\frac2r\right)+2\sqrt{kV_r},
\qquad 0\le k\lt r
\qquad\blacksquare\ \text{[Triangle Inequality; Q.E.D.]}
\end{equation*}

Finally, if $C_r=\sum_tc_t^2$ is the residue collision count, direct
expansion gives
\begin{equation*}
\begin{aligned}
V_r
&=\sum_t\left(c_t-\frac Nr\right)^2
&&[\text{By Definition}]\\
&=\sum_tc_t^2-\frac{2N}{r}\sum_tc_t+\frac{N^2}{r}
&&[\text{Expansion}]\\
&=C_r-\frac{N^2}{r}.
&&[\text{Since }\sum_tc_t=N]
\end{aligned}
\end{equation*}

Therefore a relative collision-energy theorem for the actual conditioned
starts would control the complete-block contribution to Section~5's harmful excess.
The reduction remains incomplete: no suitable relative bound for $V_r$ is
proved across the growing layers, two partial fragments remain, and the layer
bounds must still compose under the weights. When $M$ exceeds the whole
square-safe window there may be no complete block, and the boundary term
dominates.

{\raggedright
The derivation above proves the exact identities, energy bound, and
arbitrary-interval boundary. A supplementary research record is available in
\href{https://github.com/thiagomata/prime-numbers/blob/master/properties/sieve-sequence/copy-block-harmful-excess-controlled-by-residue-energy.md}{Copy-Block Harmful Excess Is Controlled By Residue Energy}.
The open relative collision input is formulated in \href{https://github.com/thiagomata/prime-numbers/blob/master/candidates/conditioned-residue-collision-energy.md}{Conditioned
Residue-Collision Energy}. The centered
rational histogram and block observable are not yet modeled as a \texttt{.holds}
theorem.
\par}
